\scanpage{6}
ou comme d'habitude l'indice $n$ indique la réduction mod \deleted{la puissance} $\underline{m}^{n+1}$ ($\underline{m}$ étant l'idéal maximal de $A$). Pour tout entier $m\quad 0$, dénotant par $ {}_mu,{}_mv$ les homomorphismes $ {}_mH\quad G$ induits par $u,v$, on aura donc aussi $({}_mu)_n=({}_mv)_n$. Ceci étant vrai pour tout $n$, et $ {}_mH$ étant \uline{fini} sur $S$, en vertu de 2.2. , il s'ensuit qu'on aura $ {}_mu={}_mv$. Ceci étant vrai pour tout $m$, on aura $u=v$ en vertu de 4.7. a).

b) Cas général. Nous \deleted{supposerons ici} pouvons supposer $S$ \deleted{\unsure{local de point fermé $s$}} \inserted{affine, $S=\operatorname{Spec}(A)$.} Quitte à faire une extension fidèlement plate et quasi compacte de la base, on peut supposer de plus $H$ diagonalisable, donc de la forme $D_S(M)$, avec $M$ un groupe \deleted{\unsure{de}}commutatif de type fini. Introduisons comme dans la démonstration de 4.6. la famille filtrante croissante des sous-$\underline{Z}$-algèbres \deleted{$A_i$} de type fini $A_i$ de $A$. Comme $\deleted{H}\unsure{G|H}$ est \unsure{ce|de} présentation finie sur $S$, il provient d'un \deleted{schéma} préschéma en groupes de présentation finie $\deleted{H}G_i$ sur un $S_i$ . D'ailleurs, $H$ provient de $H_i=D_{S_i}(M)$ De plus, pour $\deleted{i}j$ grand, en vertu de EGA IV 8, les homomorphismes \deleted{$u,v$} $u,v$ proviennent d'homomorphismes de $S_j$-groupes $\deleted{\unsure{H}}u_j,v_j:H_j\quad G_j$, \deleted{Soit $s_j$ im} et $u_j$ sera central pour $j$ grand. D'ailleurs, si $s_j$ est l'image de $s$ dans $S_j$, on a $u_{j\,s_j}=v_{j\,s_j}$ . On peut alors appliquer a) à la situation sur $S_j$, d'où aussitôt la conclusion annoncée .

\uline{Cor\deleted{\unsure{a}}ollaire} 5.2. Soient $u:H\quad G$ un homomorphisme de $S$\deleted{(}-préschémas en groupes, avec $H$ de type multiplicatif et de type fini, et $H$ de présentation finie (ou seulement localement de \deleted{pr} type fini sur $S$, lorsque $S$ est localement noethérien). Soit $s\quad S$ et supposons que $u_s$ soit l'homomorphisme unité. Alors il existe un voisinage ouvert $U$ de $s$ tel que $u_U$ soit l'homomorphisme unité.

On prouve de la même façon les résultats suivants (qui contiennent d'ailleurs les précédents, comme on voit en introduisant \deleted{\unsure{l}}\unsure{d|l}es \deleted{\unsure{grais gra|grism gra}} graphes des homomorphismes en groupes précédents):
% archive sheet number: 5

